\section{Introduction}\label{sec:introduction}

Gauge enhancement in a heterotic compactification occurs when the internal
connection acquires additional parallel adjoint sections. A deformation away
from that locus makes the corresponding external vectors massive. The
finite-dimensional Higgs mass matrix describes the leading mechanism, but
an absolute error estimate need not resolve its smallest eigenvalue: several
Higgs scales can approach zero at unrelated rates. There is a related
problem for scalar interactions. A comparison of kinetic metrics does not,
by itself, control the derivatives entering their curvature or identify the
curvature of a bundle-moduli subspace with a coupling in the full reduced
theory.

We address these questions on one family of smooth K3 compactifications of
the $\operatorname{Spin}(32)/\ZZ_2$ heterotic theory with vector structure.
We use $\SO(32)$ for its gauge algebra and its real vector-bundle
description. The background bundle contains two identical stable
rank-three factors at the enhancement point. Two orthonormal harmonic
deformation coordinates, with coefficients $a,b$, break the multiplicity
$\SU(2)$. The K3 surface remains smooth and its polarization is fixed.
The moment maps vanish throughout this bundle enhancement; no nonzero
Fayet--Iliopoulos parameter is varied.

Let $g_0$ be the fixed reference K3 metric, $V_0=\Vol(g_0)$, and let
$L^2g_0$ be the physical internal metric at leading order. The three leading
squared vector masses obey
\begin{equation}\label{eq:intro-masses}
 m_j^{2,(0)}=\widehat m_j^2\bigl(1+O(a^2+b^2)\bigr),\qquad
 (\widehat m_1^2,\widehat m_2^2,\widehat m_3^2)
 =\frac{8}{3V_0L^2}(b^2,a^2,a^2+b^2),\quad 0<b\le a.
\end{equation}
The error is proportional to each individual squared mass, with a
constant independent of $b/a$. In particular, at
$a=\varepsilon$, $b=\varepsilon^p$, $p>1$, the weakest squared mass
has order $\varepsilon^{2p}/L^2$, while its relative geometric error
still has order $\varepsilon^2$. This is a statement about actual
Hermitian--Einstein connections, obtained by applying the quantitative
local comparison of \cite{wuebben15} to the compactification family.
The Pauli-matrix eigenvalues alone would not give this conclusion.

We also calculate the first coefficient of the coupled tree-level vector
equations. Write $\eta=\alpha'/L^2$. Eliminating the complete leading
metric and two-form vector towers gives a finite, energy-dependent pencil
$F^{[1]}(\lambda,\eta)$, defined in section~\ref{sec:coupled}.
Its three positive roots, converted to squared masses by $L^{-2}$,
satisfy
\begin{equation}\label{eq:intro-corrected}
 \left|\frac{m_j^{2,[1]}}{\widehat m_j^2}-1\right|
 \le C\left(a^2+b^2+\frac{\eta}{b^2}\right).
\end{equation}
An anti-self-dual (ASD) Hodge identity checks the relative coefficients of the two
mixings and proves positivity and preservation of the unbroken kernels.
The sufficient region $\eta=o((a^2+b^2)b^2)$ makes this calculated
correction smaller than the geometric relative error. Equation
\eqref{eq:intro-corrected} concerns the defined pencil; it does not
bound an uncomputed remainder in exact string poles.

The second physical quantity is a neutral-scalar interaction. In the same
holomorphic local quotient chart, the actual harmonic metric has a
relative quadratic comparison with the hyperk\"ahler cone metric.
Ricci-flat K\"ahler regularity upgrades this comparison to every fixed
derivative order on compact smooth angular regions. An explicit
horizontal two-plane has cone sectional curvature $3/(2b^2)$.
We solve the mixed Einstein and two-form constraints, account for the
other neutral fields and for the small axion kinetic block, and obtain
the leading physical sectional curvature
\begin{equation}\label{eq:intro-curvature}
 \mathcal K_{\rm phys}^{\rm lead}
 =\frac{3g_{10}^2}{4L^2\kappa^2\varepsilon^2}
       \bigl(1+O_\kappa(\varepsilon^2)\bigr),\qquad
 a=\varepsilon,\quad b=\kappa\varepsilon,\quad 0<\kappa\le1.
\end{equation}
Here $g_{10}$ is defined by the gauge action in
section~\ref{sec:vectors}. This coefficient multiplies a specified
two-polarization quartic interaction in scalar normal coordinates.
Its derivation first extracts the tree-level large-volume coefficient
on a fixed regular patch, then takes the displayed geometric limit.
It is not uniform as $\kappa\to0$. Eliminating the Higgs multiplets
requires probe energies much smaller than $b/L$.

Heterotic bundle stability, charged fields and D-term Higgsing have an
established effective-theory description \cite{stabilitywalls}.
The $\SO(32)$ embedding, spectrum and abelian mass criteria used here
follow the K3 analysis of \cite{honecker}. Bundle and matter metrics
were computed in \cite{louis}, while the general geometry--bundle--two-form
fibration and its large-volume hierarchy were developed in
\cite{alpv}. We retain the compensator and curvature-squared terms needed
to identify the leading coefficient; we do not import a specialized
duality completion as the metric of the present compactification.
The relation between scalar curvature and covariant amplitudes is also
standard \cite{cheung}.

Universal heterotic representatives and higher-order moduli metrics are
developed in \cite{candelas,mcorist_sticka_svanes,mcorist_yin}.
Their Calabi--Yau threefold results concern four-dimensional
compactifications; we use the action and covariant-variation conventions,
not their moduli-metric theorems as results for six-dimensional K3 vacua.

The quantitative input is the energy-relative spectral and quotient-metric
comparison in \cite{wuebben15}. We use the local harmonic quadraticity
argument of \cite{wuebben20} without extending that paper's global
compactness theorem to this family. The new application combines those
inputs with a compact heterotic model, the coupled vector calculation,
and a physically normalized curvature coefficient on that same family.
The Hodge identity, interior estimates and quotient curvature formula
are standard mechanisms or direct calculations. Their role is to make
the physical identifications and their accuracy precise.

Section~\ref{sec:compactification} constructs the family and fixes its
local coordinates. Sections~\ref{sec:vectors} and~\ref{sec:coupled}
derive the vector statements and the retained field content.
Sections~\ref{sec:derivatives} and~\ref{sec:scalar} establish the
derivative estimates and the leading ambient scalar interaction.
Throughout, the string coupling is taken in the tree regime independently
of the geometric parameters.
